\begin{frame}{실데이터: diabetes 데이터셋 (2/2)}
  \begin{columns}[T]
    \begin{column}{0.44\textwidth}
      \centering
      \includegraphics[width=\linewidth]{ch02_diabetes_learning_curve.png}
      \vskip0.4em
      {\footnotesize (좌) 세 학습률의 비용 곡선, (우) 데이터 크기
      $n$ 에 따른 최종 RMSE의 학습 곡선.}
    \end{column}
    \begin{column}{0.54\textwidth}
      \small
      \textbf{(2) 학습 곡선} --- $n = 50, 100, 250, 442$ 로 학습한
      모델을 \textbf{전체 442명} 위에서 RMSE로 잴 때:
      \[
      58.7 \to 56.3 \to 53.9 \to 53.5
      \]
      데이터가 작을수록 모델이 \textbf{나빠진다} --- ``배울 것이
      적어서''.
      \vskip0.4em
      수렴 스텝 수는 $n$ 에 대해 단조가
      \emph{아님} (2869 $\to$ 912 $\to$ 2807 $\to$ 2405) ---
      ``데이터가 적으면 빠르다''는 일반 법칙이 아니라, 소음
      바닥의 크기 자체가 $n$ 에 따라 달라지기 때문.
    \end{column}
  \end{columns}
\end{frame}
